% =========================================================================
% Section III -- Problem formulation.
% The two evaluators act on the same complete schedule. Reported execution
% metrics exclude decoder failures; failure rate is reported separately.
% =========================================================================

\section{Problem Formulation}
\label{sec:problem}

\begin{table*}[t]
\caption{Notation.}
\label{tab:notation}
\centering
\scriptsize
\setlength{\tabcolsep}{3pt}
\renewcommand{\arraystretch}{1.08}
\begin{tabular}{@{}>{\raggedright\arraybackslash}p{0.13\textwidth}
                    >{\raggedright\arraybackslash}p{0.34\textwidth}
                    >{\raggedright\arraybackslash}p{0.13\textwidth}
                    >{\raggedright\arraybackslash}p{0.34\textwidth}@{}}
\toprule
Symbol & Meaning and value/domain & Symbol & Meaning and value/domain \\
\midrule
\multicolumn{2}{@{}l}{\itshape System, jobs, and schedules} &
\multicolumn{2}{l}{\itshape Evaluation and error measures} \\
$d,k,j,s,\varsigma,t,i$ & station, AMR, job, size-class, class-threshold, time/epoch, and event/instance indices &
$\widetilde\Phi,\Phi$ & idealized evaluator and collision-aware deployment executor \\
$\mathcal W$ & four-connected workspace, $\{0,\ldots,19\}\!\times\!\{1,\ldots,19\}$ &
$\pi^k,u_i^k,n_k$ & event subsequence of AMR $k$, its $i$th event, and number of events \\
$\Delta(u,v)$ & Manhattan distance between cells $u$ and $v$ &
$\ell(u),j(u)$ & station and job associated with event $u$ \\
$\mathcal Y$ & station rows, $\{2,6,10,14,18\}$ &
$\widetilde C_k,\widetilde C$ & idealized completion time of AMR $k$ and idealized makespan \\
$\mathcal D^{\rm in},\mathcal D^{\rm out},\mathcal D$ & inbound stations, outbound stations, and their union &
$H$ & space--time A* search horizon, $100$ ticks \\
$\bm u_d,W_d,L$ & inward wall normal, waiting cells of station $d$, and queue length $L=3$ &
$\nu_\Phi(x)$ & routing or operation events that $\Phi$ cannot decode \\
$\mathcal M,m$ & AMR set and fleet size, $m=16$ &
$\mathcal F_\Phi$ & schedules routable by the fixed executor, i.e., $\nu_\Phi(x)=0$ \\
$\mathcal B,b_k$ & set of dedicated bays and the bay assigned to AMR $k$ &
$F_k^\Phi,C_{\max}^\Phi$ & executed completion time of AMR $k$ and executed makespan \\
$\mathcal J,n$ & job set and number of jobs &
$x^*$ & executor-relative makespan-minimizing schedule \\
$\mathcal S,\sigma_j$ & size classes $\{A,B,C\}$, $A\prec B\prec C$, and class of job $j$ &
$\Lambda(x)$ & relative value error (execution penalty) of $\widetilde\Phi$ \\
$o_j,g_j$ & inbound origin and outbound destination of job $j$ &
$\mathcal C_i$ & common executor-routable candidate set for test instance $i$ \\
$p_j=p(\sigma_j)$ & pickup/delivery service time; $p(A),p(B),p(C)=5,10,15$ &
$P_i,D_i$ & concordant and discordant schedule-pair counts \\
$\rho_j,\delta_j$ & pickup and delivery events of job $j$ &
$T_i^{\widetilde\Phi},T_i^\Phi$ & pair counts tied only by the indicated evaluator \\
$Q_s,N_s^k(t)$ & class-$s$ rack positions ($Q_s=1$) and onboard class-$s$ jobs &
$\tau_{b,i},{\rm Inv}_i$ & tie-corrected Kendall coefficient and inversion rate \\
$x=(\bm a,\pi)$ & complete schedule: assignment vector and global event-priority list &
$\widetilde x_i,x_i^\Phi,{\rm Reg}_i$ & selections of the two evaluators and executor-relative selection regret \\
$\bm a,\pi$ & job-to-AMR assignments and a precedence-feasible list of $2n$ events &
& \\
$\mathcal X$ & schedules satisfying assignment, precedence, and rack constraints &
& \\
\multicolumn{2}{@{}l}{\itshape Constructive decision process} &
& \\
$s_t$ & state at construction epoch $t$ & & \\
$\bm X_t^A,\bm X_t^J$ & projected AMR and job states & & \\
$\bm E_t,\pi_{1:t}$ & committed station-service windows and current event-list prefix & & \\
$a_t=(\omega,k,j)$ & composite action selecting operation, AMR, and job & & \\
$\omega$ & operation type, $\textsc{pick}$ or $\textsc{del}$ & & \\
$\widehat\Psi$ & calibrated event-based surrogate transition & & \\
\bottomrule
\end{tabular}
\end{table*}

\subsection{System, Jobs, and Schedules}
\label{sec:system}

Table~\ref{tab:notation} provides a consolidated reference for the notation
used in this section. The workspace is the obstacle-free, four-connected grid
$\mathcal W=\{0,\ldots,19\}\times\{1,\ldots,19\}$ with unit-time moves and
Manhattan distance $\Delta$. Five inbound stations and five outbound stations
occupy the left and right boundaries at rows
$\mathcal Y=\{2,6,10,14,18\}$:
\begin{equation}
 \mathcal D^{\mathrm{in}}=\{(0,y):y\in\mathcal Y\},\qquad
 \mathcal D^{\mathrm{out}}=\{(19,y):y\in\mathcal Y\}.
 \label{eq:stations}
\end{equation}
Let $\mathcal D=\mathcal D^{\mathrm{in}}\cup\mathcal D^{\mathrm{out}}$.
Each station is an exclusive, non-preemptive server. For a station $d$, let
$\bm u_d$ be the inward wall normal. Its finite single-file waiting set is
\begin{equation}
 W_d=\{d+r\bm u_d\in\mathcal W\setminus\mathcal D:
 r\in\{1,\ldots,L\}\},\qquad L=3.
 \label{eq:waitset}
\end{equation}
Thus, each queue extends three cells inward from its station and leaves the
two lateral departure cells free; each waiting cell holds at most one AMR.
The fleet is fixed at $m=16$. Its dedicated bays occupy four aisle rows
between station rows,
\begin{equation}
 \mathcal B=\{0,1,2,3\}\times\{4,8,12,16\},
 \label{eq:bays}
\end{equation}
with a one-to-one assignment $b_k\in\mathcal B$ for each AMR
$k\in\mathcal M$. An idle AMR occupies its bay, and every AMR returns to its
bay after completing its assigned work. Figure~\ref{fig:layout} shows this
fixed geometry.

An instance contains jobs $\mathcal J=\{1,\ldots,n\}$, where each job
represents one parcel transfer. Job $j$ has size
$\sigma_j\in\mathcal S=\{A,B,C\}$, with $A\prec B\prec C$, origin
$o_j\in\mathcal D^{\mathrm{in}}$, destination
$g_j\in\mathcal D^{\mathrm{out}}$, and deterministic service time
\begin{equation}
 p_j=p(\sigma_j),\qquad p(A)=5,\quad p(B)=10,\quad p(C)=15.
 \label{eq:service}
\end{equation}
The same duration is required for pickup $\rho_j$ and delivery $\delta_j$,
and all jobs are available at time zero. Pickup and delivery must be
performed by the same AMR; transshipment is not allowed.

Each AMR has one rack position of each class. A smaller job may occupy a
larger position, but not conversely. With $N_s^k(t)$ denoting the number of
onboard class-$s$ jobs and $Q_s=1$ for every $s\in\mathcal S$, the nested
capacity constraints are
\begin{equation}
 \sum_{s\succeq\varsigma}N_s^k(t)
 \le \sum_{s\succeq\varsigma}Q_s,\qquad
 k\in\mathcal M,\ \varsigma\in\mathcal S,\ t\ge0.
 \label{eq:nested}
\end{equation}
Equivalently, $N_C^k\le1$, $N_B^k+N_C^k\le2$, and
$N_A^k+N_B^k+N_C^k\le3$.

A complete schedule is $x=(\bm a,\pi)$. The vector
$\bm a\in\mathcal M^n$ assigns each job to one AMR, while $\pi$ is a priority
list containing the $2n$ pickup and delivery events. It obeys
$\rho_j\prec_\pi\delta_j$, and every prefix of each robot's induced
subsequence satisfies~\eqref{eq:nested}. Let $\mathcal X$ denote the set of
schedules satisfying these assignment, precedence, and rack constraints.
The list $\pi$ fixes event priority, not physical start times: station service,
waiting, and routed movement are supplied by an evaluator.

\subsection{Paired Schedule Evaluators}
\label{sec:evaluators}

We apply two deterministic evaluators to the same $x\in\mathcal X$. For
robot $k$, let $n_k$ be its number of events and
$\pi^k=(u_1^k,\ldots,u_{n_k}^k)$ the subsequence induced by $x$. Let
$\ell(u)$ and $j(u)$ be an event's station and job, respectively, and set
$\ell(u_0^k)=b_k$. The idealized evaluator $\widetilde\Phi$ assumes
independent robots, immediate service on arrival, and free-space travel:
\begin{equation}
 \begin{aligned}
 \widetilde C_k(x)
 &=\sum_{i=1}^{n_k}\!\left[
   \Delta\!\left(\ell(u_{i-1}^k),\ell(u_i^k)\right)+p_{j(u_i^k)}\right]\\
 &\quad+\Delta\!\left(\ell(u_{n_k}^k),b_k\right),
 \qquad
 \widetilde C(x)=\max_{k\in\mathcal M}\widetilde C_k(x).
 \end{aligned}
 \label{eq:ideal}
\end{equation}
It therefore ignores inter-robot conflicts, station exclusivity, finite
waiting space, and delays caused by reservations. Global interleavings with
the same per-robot subsequences may consequently tie under
$\widetilde\Phi$.

The deployment evaluator $\Phi$ is a fixed prioritized executor. It processes
events in $\pi$, routes each leg with space--time A* over a shared reservation
table, and admits unit moves or waits. A move is rejected if its next cell is
reserved at the current or next tick; this prevents vertex conflicts,
head-on exchanges, and immediate following into a vacated cell. Each station
serves one AMR at a time. If service cannot start on arrival, the executor
reserves an available cell in $W_d$; when none is available, the AMR remains
at its upstream location until the next relevant release time. After its
last event, each AMR is routed back to $b_k$.

The space--time search horizon is $H=100$ ticks. Let $\nu_\Phi(x)$ count
routing or operation events that the decoder cannot complete within its
rules, and define the executor-relative routable set
\begin{equation}
 \mathcal F_\Phi=\{x\in\mathcal X:\nu_\Phi(x)=0\}.
 \label{eq:routable-set}
\end{equation}
For $x\in\mathcal F_\Phi$, let $F_k^\Phi(x)$ include all movement, waiting,
service, and the return to $b_k$. The executed makespan and optimization
objective are
\begin{equation}
 C_{\max}^{\Phi}(x)=\max_{k\in\mathcal M}F_k^\Phi(x),
 \qquad
 x^*\in\arg\min_{x\in\mathcal F_\Phi}C_{\max}^{\Phi}(x).
 \label{eq:objective}
\end{equation}
Because $\Phi$ is a bounded prioritized planner, $\mathcal F_\Phi$ is
executor-relative rather than a claim of general MAPF feasibility. A failed
decode contains a fixed horizon penalty rather than elapsed execution time.
We therefore report routability separately and compute all execution-time and
ranking statistics below only on schedules in $\mathcal F_\Phi$.

\subsection{Value and Ranking Error}
\label{sec:penalty}

For the same executor-routable schedule, the relative value error of the
idealized evaluator is the execution penalty
\begin{equation}
 \Lambda(x)=
 \frac{C_{\max}^{\Phi}(x)-\widetilde C(x)}{\widetilde C(x)}.
 \label{eq:penalty}
\end{equation}
Value error alone does not show whether the abstraction changes a scheduling
decision. For test instance $i$, let $\mathcal C_i\subseteq\mathcal F_\Phi$
be the common set of complete schedules produced by all evaluated methods.
For every unordered pair $\{x,y\}\subset\mathcal C_i$, define it as concordant
when
\begin{equation}
 [\widetilde C(x)-\widetilde C(y)]
 [C_{\max}^{\Phi}(x)-C_{\max}^{\Phi}(y)]>0,
 \label{eq:concordant}
\end{equation}
and discordant when the product is negative. Let $P_i$ and $D_i$ be the
respective counts, and let $T_i^{\widetilde\Phi}$ and $T_i^\Phi$ count pairs
tied only by the indicated evaluator. Kendall's tie-corrected coefficient is
\begin{equation}
 \tau_{b,i}=
 \frac{P_i-D_i}
 {\sqrt{(P_i+D_i+T_i^{\widetilde\Phi})
              (P_i+D_i+T_i^\Phi)}}.
 \label{eq:kendall}
\end{equation}
The directly measured ranking-inversion rate is
\begin{equation}
 \mathrm{Inv}_i=\frac{D_i}{P_i+D_i},
 \label{eq:inversion}
\end{equation}
so pairs tied by either evaluator are excluded from its denominator and are
reported separately. Finally, using a predetermined tie-breaking rule, let
$\widetilde x_i=\arg\min_{x\in\mathcal C_i}\widetilde C(x)$ and
$x_i^\Phi=\arg\min_{x\in\mathcal C_i}C_{\max}^{\Phi}(x)$. Their
executor-relative selection regret is
\begin{equation}
 \mathrm{Reg}_i=
 \frac{C_{\max}^{\Phi}(\widetilde x_i)-C_{\max}^{\Phi}(x_i^\Phi)}
 {C_{\max}^{\Phi}(x_i^\Phi)}.
 \label{eq:selection-regret}
\end{equation}
The three ranking measures are computed per instance before aggregation, so
instances with more sampled schedules do not receive unintended extra weight.

\subsection{Constructive Decision Process}
\label{sec:mdp}

The policy constructs one schedule in $2n$ epochs. At epoch $t$, its state and
composite action are
\begin{align}
 s_t&=(\bm X_t^A,\bm X_t^J,\bm E_t,\pi_{1:t}),
                                                        \label{eq:state}\\
 a_t&=(\omega,k,j),\qquad
 \omega\in\{\textsc{pick},\textsc{del}\}.             \label{eq:action}
\end{align}
Here $\bm X_t^A$ and $\bm X_t^J$ describe the projected robot and job states,
$\bm E_t$ stores committed station service windows, and $\pi_{1:t}$ is the
current prefix. The action mask allows pickup of an unpicked job only by an
AMR with compatible rack capacity, and allows delivery only after pickup by
that job's carrier. It therefore enforces assignment, precedence, and nested
capacity during construction.

Both training and inference advance a prefix with the calibrated event-based
surrogate
\begin{equation}
 s_{t+1}=\widehat\Psi(s_t,a_t),
 \label{eq:surrogate-transition}
\end{equation}
which supplies inexpensive projected travel, station, and waiting states. It
does not determine the optimization target. Once a complete schedule has
been constructed, the fixed executor $\Phi$ evaluates it; PPO receives the
negative executor-evaluated makespan as its terminal return on routable
rollouts. Failed rollouts receive the executor's finite training-only score,
while their returned duration is excluded from reported makespan statistics.
Section~\ref{sec:architecture} specifies $\widehat\Psi$, the observable
congestion features, and the policy parameterization.

\subsection{Scope and Assumptions}
\label{sec:assumptions}

Robots are homogeneous, all jobs are known before construction, and service
times are deterministic and class-dependent. Here \emph{offline} describes
the scheduling instance, not offline reinforcement learning. The controlled
study fixes the layout, the $m=16$ fleet, and the executor, while varying
workload through $n$. Battery consumption, turning and acceleration costs,
localization error, transshipment, and online arrivals are excluded.
Consequently, the results identify workload-dependent evaluator error within
this controlled system; generalization to other layouts, fleets, and routing
executors remains outside the present scope.
